% Generated draft: requires mathematical and source review.
\documentclass[11pt]{article}
\usepackage[margin=25mm]{geometry}
\usepackage{fontspec}
\setmainfont{Latin Modern Roman}
\usepackage{amsmath,amssymb,mathtools,unicode-math}
\setmathfont{Latin Modern Math}
\usepackage{enumitem,needspace,xurl,hyperref}
\setlength{\parindent}{0pt}
\setlength{\parskip}{5pt}
\setlength{\emergencystretch}{3em}
\newcommand{\N}{\mathbb{N}}
\newcommand{\Z}{\mathbb{Z}}
\newcommand{\Q}{\mathbb{Q}}
\newcommand{\R}{\mathbb{R}}
\newcommand{\C}{\mathbb{C}}
\newcommand{\F}{\mathbb{F}}
\newcommand{\A}{\mathbb{A}}
\newcommand{\PP}{\mathbb{P}}
\newcommand{\E}{\mathbb{E}}
\newcommand{\eps}{\varepsilon}
\newcommand{\dd}{\,\mathrm{d}}
\begin{document}
\section*{UnsolvedMath ID: 30003823}
\textbf{Source problem code:} catalogue entry
\subsection{Definitions and mathematical statement}
The supplied problem asks whether, for every finite reflection group \(W\), its pure braid group admits the presentation
\[
P_c(W)\stackrel{?}{\cong}\langle T:\operatorname{Red}_T(\Delta)\rangle,
\]
where \(T\) consists of squares of reflection generators and \(\Delta\) is the full twist. The source does not specify the class of reflection groups, the meaning of the subscript \(c\), the reflection generators and their braid lifts, or the relations encoded by \(\operatorname{Red}_T(\Delta)\). These choices must be resolved before this is a uniquely determined mathematical assertion.

One possible precise interpretation of the presentation notation is the following, conditional on the intended choices. Write \(P=P_c(W)\), and suppose that specified braid lifts \(\widetilde r\) of specified reflection generators give elements
\[
T=\{\widetilde r^{\,2}\}\subseteq P.
\]
Let \(\Delta\in P\) be the specified full twist. Assume that \(\Delta\) has a positive factorization over \(T\), and define
\[
\ell_T^+(\Delta)=\min\{k\geq 0:\Delta=t_1\cdots t_k,\ t_i\in T\}.
\]
Under the interpretation that reduction means shortest positive factorization, put
\[
\mathcal R_T(\Delta)=\{(t_1,\ldots,t_k):k=\ell_T^+(\Delta),\ t_i\in T,\ t_1\cdots t_k=\Delta\text{ in }P\}.
\]
Introduce a formal generator \(x_t\) for each \(t\in T\), and define
\[
G(T,\Delta)=\left\langle x_t\ (t\in T)\ \middle|\ x_{t_1}\cdots x_{t_k}=x_{u_1}\cdots x_{u_k}\ \text{for all }(t_1,\ldots,t_k),(u_1,\ldots,u_k)\in\mathcal R_T(\Delta)\right\rangle.
\]
The evaluation map \(x_t\mapsto t\) defines a homomorphism \(G(T,\Delta)\to P\). On this interpretation, the question is whether this homomorphism is an isomorphism for every \(W\) in the intended class, with the intended choices. This asks both whether \(T\) generates \(P\) and whether the displayed relations give all relations among those generators.

This conditional expansion is not asserted to be the source's definition of \(\operatorname{Red}_T(\Delta)\). In particular, positive factorization, the choice of lifts, and the pairwise equality interpretation of the relations require confirmation. The supplied unexpanded question remains the target until these qualifications are resolved.
\subsection{Short English statement}
The question asks whether every intended finite reflection group has a pure braid group presented by squared reflection generators, with all defining relations coming from reduced factorizations of the full twist. The source leaves essential choices and the relation notation unspecified, so the explicit presentation above is conditional. The supplied record labels the universal question open; that label is not independently verified here.
\subsection{Sources}
\begin{itemize}
\item Local UnsolvedMath registry, supplied complete record for permanent ID \(30003823\), titled \emph{Uniform Presentations of Pure Reflection Braid Groups}; source problem code \texttt{catalogue entry}, in the Oberwolfach Reports Open Problems collection. The record gives proposed year 2018 and names no individual proposer.
\item Supplied source citation: \emph{Enumerative Combinatorics (2019)}, Oberwolfach Reports, DOI 10.4171/owr/2018/23, pp. 1381--1464. Original OWR report: \url{https://doi.org/10.4171/owr/2018/23}. The supplied background describes Oberwolfach Reports as documenting mathematical workshops and their open-problem discussions.
\item Supplied literature source: Liu--Williams full-twist presentations, \url{https://math.mit.edu/research/highschool/primes/materials/2025/Liu-Williams.pdf}.
\item Literature assessment dated exactly \texttt{2026-08-21}: the supplied curation labels the problem open and reports proofs in Coxeter types \(A\), \(B\), \(D\), \(H_3\), \(I_2(m)\), and \(F_4\), with the remaining exceptional types still open. This is supplied dated curation, not an independently verified literature assessment. The background explicitly calls it dated literature triage derived from the supplied OWR curation and requests independent verification of time-sensitive status and cited claims.
\end{itemize}
\end{document}
